\section{Uniform dependence on the detuning bound}\label{inf:sec-tame}

Fix the compact split interval, a finite centre accuracy $N$, and
every finite derivative order and complex disc required for that
centre. Constants in this section depend on these choices.
Coefficient and mixed operator estimates are at integer block sizes;
the intervening flow parameters use the weighted realization of
Proposition~\ref{inf:flow-form}.
They are independent of $p$ and of $0<\tau\le T\le\log p$.
We may replace $T$ by $\max(1,T)$ when forming an envelope.

\begin{lemma}[Detuning-dependent estimates]\label{inf:detuning-dependencies}
Suppose $|p-p_0|\le C_{\mathrm{ap}}/p_0$, where
$C_{\mathrm{ap}}\le C_*\log p_0$ with fixed $C_*$, and the continued
zero has $0<\tau\le T\le\log p$. There are finite integers
$k_N,d_N^{\mathrm{tame}}$ and constants $c_N,C_N>0$ such that the
finite-centre estimates hold on
\[
 |\eps|\le c_N(1+T)^{-k_N},\qquad
 \mathfrak P_N(T)=(1+T)^{d_N^{\mathrm{tame}}},\qquad
 \mathcal B_N(T)=C_N\mathfrak P_N(T)e^{C_N\sqrt T}.
\]
Increasing the fixed exponents and constants when necessary, every
bound in the following table holds. Here $E=K_{\mathrm{flow}}(1+\tau)$,
$W=C_{\mathrm{cond}}(1+\tau)$, and $\delta$ is the strong centre size.
The labels in the first column specify the inherited statement
and the quantity whose bound is being extended.

\begingroup\small\normalfont
\renewcommand{\arraystretch}{1.16}
\begin{longtable}{@{}>{\raggedright\arraybackslash}p{.06\textwidth}>{\raggedright\arraybackslash}p{.38\textwidth}>{\raggedright\arraybackslash}p{.48\textwidth}@{}}
\toprule & Quantity and statement & Uniform bound / small parameter\\\midrule
\endfirsthead
\toprule & Quantity and statement & Uniform bound / small parameter\\\midrule
\endhead
1 & Root continuation, \ref{inf:detuning} &
$R-R_{\mathrm{alg}}=-(2\delta_0+O(p^{-1}))(p-p_0)$.\\
2 & Denominator in \eqref{inf:trace-ratio-exact} &
$2R^2-4(p+1)(p+2)=4p^2(1+O(p^{-1}))$.\\
3 & Low-state transfer, \ref{inf:wide-low-transfer} &
$C C_{\mathrm{ap}}(\ell+1)/p+C C_{\mathrm{ap}}(1+C_{\mathrm{ap}})/p^2$.\\
4 & Radius equation, \eqref{inf:rho-equation} &
Analytic in $(\eps,w)$, $w=\theta\eps^2$; $|w|\le T/p^2$.\\
5 & Finite centre jet, \ref{inf:centres} &
Each fixed $\eps$ coefficient is polynomial in $\theta$;
common $\eps$ radius $c_N(1+T)^{-k_N}$.\\
6 & Collar and complex angular disc, \eqref{inf:collar-ode} &
Fixed discs; physical displacement $C\tau/p$.\\
7 & Two physical defects, \eqref{inf:centre-complex-defect} &
$C_N\theta^2p^{-N}\mathfrak P_N(T)$, also in the weighted
coefficient norms of \ref{inf:centre-coefficient-defect}.\\
8 & Quartic comparison, \eqref{inf:centre-profile} &
Physical error $C_N\tau p^{-3}\mathfrak P_N(T)$;
normalized error $C_N\tau p^{-4}\mathfrak P_N(T)$.\\
9 & Density, \eqref{inf:density-graph-bounds} &
$\nrm{m^{\pm1}}\le e^{Cp\delta}$,
$\delta\le C\tau/p^2$, so $p\delta\le CT/p$.\\
10 & Radial modes, \ref{inf:radial-centre-norm} &
$C_{N,k}p^{k+1}$; normalization units remain separated from zero.\\
11 & Known-mode composition, \eqref{inf:composition-bound} &
Head factor $e^{C_N\sqrt T}$; tail step factor
$e^{C_N\sqrt T/p}$ against Bessel ratio $1/4$.\\
12 & Clamped centre and material,
\ref{inf:centre-source-norm}, \ref{inf:sharp-material} &
$\nrm{U_0}\le p\mathcal B_N(T)$;
$\nrm{g_0}\le p^3\mathcal B_N(T)$;
$\nrm{\Qmat^{-1}}\le p^2\mathcal B_N(T)$;
$\sup\nrm{D^2\Fnl}\le p^3\mathcal B_N(T)$.\\
13 & Relative flow form, \ref{inf:flow-relative} &
$C\tau p^{-2}$ and $C\tau^2p^{-3}$;
series ratio $C\tau/p^2$.\\
14 & Flow complement, \ref{inf:flow-complement} &
Choose $E=K_{\mathrm{flow}}(1+\tau)$; ratio $C\tau/E$.\\
15 & Schur derivative, \eqref{inf:flow-schur-remainder} &
$C\tau^2(E^{-1}+p^{-1})$;
$\tau^3/E^2\le\tau^2/K_{\mathrm{flow}}E$.\\
16 & Seed diagonal, \ref{inf:flow-low-block} &
$\theta/(2p)+O(\theta^2p^{-3})$;
weighted low remainder $C\tau^2(E^{-1}+p^{-1})p^{-2}$.\\
17 & Index-five diagonal, \eqref{inf:low-offsets} &
$-72+O((1+T)/p)$; primary perturbation $O(\tau/p)$.\\
18 & Low feedback, \eqref{inf:weighted-low-inverse} &
For $a_*=C\tau^2(E^{-1}+p^{-1})$, bound $Ca_*^2/(p\tau)$;
ratio to $\tau$ tends to zero.\\
19 & Transport, \eqref{inf:transport-bounds} &
$C\tau^2/p$ in the weighted derivative form;
ratio to the main margin $O(\tau/p)$.\\
20 & Norm recovery, \ref{inf:flow-theorem} &
$\nrm u\le C(\sqrt p+\tau)\nrm{\mathsf Wv}$;
slope at least $c\tau/p$.\\
21 & Phase/count, \ref{inf:window-isolation}--\ref{inf:global-count} &
Phase thickness $C(1+W)/p$;
count $C\sqrt{1+W}\,p^{2/3}$ for $W\le C\log p$.\\
22 & Shell and split count, \ref{inf:quartic-count-window},
\ref{inf:many-splits} &
Shell $C\tau/p^2$; count window $C(1+\tau)$;
bad splits $C\sqrt{1+\tau}(p^{2/3}+p^{8/3-q_{\mathrm{gap}}})$.\\
23 & Radial conditioning, \ref{inf:radial-poles} &
Complementary graph $Cp$, Neumann ratio $C\tau/p$.\\
24 & Linear spectral comparison, \ref{inf:centre-form-comparison} &
Error $C\tau^2p^{-2}$; divided by $\gamma$ it is
$C\tau p^{q_{\mathrm{gap}}-2}$.\\
25 & Angular conditioning window &
$W=C_{\mathrm{cond}}(1+\tau)$; every principal inverse $2/W$.\\
26 & High core, \ref{inf:critical-strips} &
Diameter $L_{\mathrm{core}}+C\sqrt{W+1}$;
condition factor $e^{C(L_{\mathrm{core}}+\sqrt{W+1})}
\sqrt{L_{\mathrm{core}}+C\sqrt{W+1}+1}$.\\
27 & High barrier, \ref{inf:off-core-table} &
$\log(L_{\mathrm{core}}^2/[C(1+\tau)])\ge
\tfrac12\log\log p$; choose $K_{\mathrm{core}}$ before $p$.\\
28 & Low core, \ref{inf:path-valuation} &
$j\le C\sqrt{1+\tau}$; one-step coefficient
$C(j+1)/p$, downward coefficients $C(j+1)^2/p^2$.\\
29 & Low-core off-diagonal, \ref{inf:off-core-table} &
$C\tau(1+\tau)^2p^{-2}$ in coefficient and Hilbert norms.\\
30 & Singleton off-core, \eqref{inf:global-off-core} &
$C\tau(1+\tau)^2p^{-2+2\alpha_{\mathrm{loc}}}$.\\
31 & True-centre graph, \ref{inf:quadratic-graph} &
$C_N\tau p^{-4}(\mathfrak P_N(T)+\tau)$;
restore in $\Mmix$ before coefficient conversion.\\
32 & Radial conversion, \ref{inf:head-tail} &
One factor $\sqrt p$; tail inverse uniformly bounded.\\
33 & Convexity, \ref{inf:geometry-convexity} &
Normalized $C^2$ shape $C\tau p^{-2}$ plus Newton correction.\\
34 & Non-ball witness, \ref{inf:geometry-nonball} &
Leading physical coefficient $-\theta/(16r(1-r)p)$;
centre error $C_N\tau p^{-2}\mathfrak P_N(T)$.\\
\bottomrule
\end{longtable}
\endgroup
The double lift, residual, and material-defect bounds are
\begin{equation}\label{inf:tame-residuals}
 \nrm{\ell_N}\le C_N\theta^2p^{-N}\mathfrak P_N(T),\quad
 \nrm{\Fnl(z_0)}\le C_N\theta^2p^{2-N}\mathfrak P_N(T),\quad
 \nrm{\mathcal V_{\mathrm{res}}}\le
                  C_N\theta^2p^{3-N}\mathfrak P_N(T).
\end{equation}
The normal Hessian satisfies
$\nrm{D_p^{-1}}_{\Ec_1}\le C_Np^{-2}\mathcal B_N(T)$.
The scalar fixed-centre graph difference is $C\tau p^{-2}$ and
the unitary relative bound is $C(\delta+p\delta^2)$.
Every $o(1)$ obtained from these envelopes is uniform in
$0<\tau\le\log p$. The local cluster/cofactor bound is used only
with fixed $T$; the growing-window inverse is
Theorem~\ref{inf:mixed-conditioning}.
\end{lemma}
\begin{proof}
We give the dependence calculations in the order of their inputs.

\emph{Roots and low states.}
Equation~\eqref{inf:order-derivative} has error $O(p^{-2})$ and its
comparison function is smooth near $1/2$. Throughout the stated
interval, $p/R=1/2+O(p^{-1})$ and $R\to\infty$.
Thus $R'=2\pi/(3\sqrt3)+O(p^{-1})$ and
$R'_{\mathrm{alg}}=2+O(p^{-2})$.
Integrating their difference retains exactly the factor $p-p_0$.
Substitution in \eqref{inf:trace-ratio-exact} proves rows 1--2 and
\begin{equation}\label{inf:log-detuning-relation}
 \theta=(16\delta_0+O(p^{-1}))p_0(p-p_0),\qquad
 R=2p+3+O((1+C_{\mathrm{ap}})/p).
\end{equation}
There is no additive detuning error. Row 3 is
\eqref{inf:wide-low-error}.

\emph{Finite jets and their remainders.}
The radius equation is analytic in $(\eps,w)$ with
$w=\theta\eps^2$. Coefficient extraction in the finite recurrence,
the collar ODE, the monic products, and the update
\eqref{inf:formal-update} shows inductively that every fixed
$\eps$ coefficient is a polynomial in $\theta$.
All leading divisors are independent of $\theta$: they are the
normalization units, $-1/8$, $c_j$, and nonzero functions of $r$.
For this fixed finite jet choose $k_N$ larger than the finitely
many ratios of positive $\theta$ degrees to positive $\eps$ orders.
On $|\eps|\le c_N(1+T)^{-k_N}$, $|\theta|\le2(T+1)$, its
amplitudes and displacement are uniformly small by choosing $c_N$.
The variable $w$, the collar denominators $\rho+\eps y$, and every
finite normalization remain in fixed compact analytic sets.
The same choice works on the prescribed fixed complex angular disc.

The exact defects of the evaluated finite jet vanish through the
required order in $\eps$. They have an analytic factor $\theta^2$
by the identical field/boundary linear seed, as proved in
Theorem~\ref{inf:centres}. Divide by this factor; the maximum
principle on the larger $\theta$ circle bounds the quotient
uniformly on the smaller disc. The Cauchy remainder estimate in
$\eps$ costs at most a fixed power $(1+T)^{d_N^{\mathrm{tame}}}$.
The same argument applied to $(h_N-h_{\mathrm q})/\theta$ uses its
vanishing first three $\eps$ coefficients and gives row 8.
For $T\le\log p$ the condition
$p^{-1}(1+T)^{k_N}\to0$ puts the real input inside the chosen disc.
Lemma~\ref{inf:analytic-coefficients} transfers the bounds to every
fixed coefficient moment and radius. This proves rows 4--8.

The first term of the normalized shape has size $C\tau p^{-2}$,
with a constant independent of the fixed order. Its remainder is
$C_N\tau p^{-4}\mathfrak P_N(T)$; its ratio to the leading size
tends to zero. Thus $\delta\le C\tau p^{-2}$ in the strong norm.
The scalar estimates in the proof of Theorem~\ref{inf:relative-graph}
are $C(\delta+p\delta^2)$, with no fixed-$T$ compactness constant.
Since $p\delta=O(T/p)\to0$, the density series give row 9 and the
stated graph bounds.

The mode-normalization ratios in
Lemma~\ref{inf:radial-centre-norm} are finite recurrence functions
of $(\eps,w)$ and tend to the same nonzero constants. Its Parseval
head bound and the real-order tail ratio therefore hold uniformly.
For composition, the argument of the square root in
\eqref{inf:composition-bound} is $C_N\tau p^{-2}s(2s+\beta+1)$.
It is $O_N(T)$ on the $s=O_N(p)$ head, and on the tail gives an
extra geometric step $e^{C_N\sqrt T/p}$, absorbed by the ratio
$1/4$. The finite amplitudes stay bounded on the real input.
This proves rows 10--11 and the old membership/moment estimates
with envelope $\mathcal B_N(T)$.
The exact lift columns and their Laplacians give
\eqref{inf:tame-residuals} by precisely the two-layer calculation
of Proposition~\ref{inf:residual-estimates}.
The inverse multiplier series preserve every fixed moment since
$C\tau/p^2\to0$.
The normal Hessian identity \eqref{inf:normal-hessian} now gives
its stated inverse bound. Repeating the two norm absorptions in
Proposition~\ref{inf:sharp-material} yields row 12; finite products
of its envelopes are included by increasing $d_N^{\mathrm{tame}},C_N$.
These arguments use composition only for the known finite modes,
clamped gains only on $\Gc_p$, and all-layer Poisson estimates
for $\mathfrak f_0$.

\emph{Uniform flow.}
The series in Lemma~\ref{inf:flow-relative} is dominated by
$\sum h^6(C\tau/p^2)^h$, proving row 13 with an absolute constant
for the fixed split interval. In the complementary proof choose
$E=K_{\mathrm{flow}}(1+\tau)$ independently of $\varsigma$.
Then $C\tau/E\le C/K_{\mathrm{flow}}$ and $E/p^2\to0$.
The differentiated complementary inverse adds
$C\tau^3/E^2$, bounded as in row 15. Thus rows 14--15 hold
without differentiating a moving window projection.
The finite recurrence and the collar root argument give rows
16--17 with their displayed factors of $\theta$.
In the low feedback calculation
\[
 \frac{a_*^2}{p\tau}
 \le C\left(\frac{\tau}{pK_{\mathrm{flow}}^2}
       +\frac{\tau^2}{p^2K_{\mathrm{flow}}}
       +\frac{\tau^3}{p^3}\right)=o(\tau).
\]
The transport error is $C\tau^2/p=o(\tau)$.
Choose $K_{\mathrm{flow}}$ to absorb $C\tau/K_{\mathrm{flow}}$
into the positive angular margin and then take $p$ large.
Low-state transfer from row 3 gives the same $32$ exclusion on
indices below $\lfloor\sqrt p/16\rfloor$.
The norm recovery in the last paragraph of
Theorem~\ref{inf:flow-theorem} is
$\nrm u\le C(\sqrt p+\tau)\nrm{\mathsf Wv}$;
$\tau\le\log p=o(\sqrt p)$ proves rows 18--20.

\emph{Counting.}
The proof of Lemma~\ref{inf:window-isolation} needs only
$2W/R<1/10$ and a frequency interval shorter than $\pi$.
These hold for $W\le C\log p$.
In Lemma~\ref{inf:phase-strip}, the Debye error is independent of
$W$ and replacing the zero by $R$ costs at most $C W/p$.
Its thickness is therefore $C(1+W)/p$.
On each dyadic band the determinant error in
Lemma~\ref{inf:lattice-strip} has its same fixed curvature term
and additional error $O(Wp^{-2/3})=o(1)$.
Its number of collinear points is at most $C\sqrt{1+W}$.
Summing the dyadic intervals as in Theorem~\ref{inf:global-count}
gives row 21; the turning strip has only $O(p^{2/3})$ indices
uniformly in this range of $W$.
The shell width $C\tau/p^2$ gives count window $C(1+\tau)$ by
indexwise min--max. The grid proof of Lemma~\ref{inf:many-splits}
then gives row 22 with its explicit square-root factor.

\emph{Conditioning and geometry.}
Rows 23--32 are the estimates proved in
Section~\ref{inf:sec-conditioning}, which kept $\tau\le\log p$
throughout. Row 8 supplies its true-centre hypothesis with
$\mathfrak E_N(p)=C_N\mathfrak P_N(\log p)$.
The high-core diameter in row 26 is linear in $\sqrt W$;
its exponential is therefore subpolynomial. Its barrier gain is
row 27. The three low-core facts used in rows 28--29 are: its
indices have residue two modulo three; every primary edge lands
on a nonresonant index with denominator $cp$; a feedback edge
that lands on another resonant index has displacement three and
the extra coefficient $(j+1)/p$. These are the statements proved
in Lemmas~\ref{inf:critical-strips} and \ref{inf:path-valuation}.
The full off-diagonal estimates are operator norms as in
Lemma~\ref{inf:off-core-table}. The factor in row 31 is multiplied
by the mixed graph inverse, yielding
$p^{q_{\mathrm{gap}}+\xi_{\mathrm{inv}}-2+o(1)}\to0$.
The final radial conversion costs just row 32.
Rows 33--34 now follow from the fixed-order shape comparisons and
the dimension-independent $C^2$ embedding.

Finally
\[
 \mathcal B_N(\log p)
 =\exp\{O_N(\log\log p+\sqrt{\log p})\}=p^{o(1)}.
\]
The other envelopes in the table are fixed powers of $\log p$ or
have the form
\[
 \exp\{O(\log p/\log\log p+\sqrt{\log p})\}.
\]
The strict power margins absorb all of them. The two factors in
\eqref{inf:tame-residuals} are $\theta^2$, the inverse has only
$\tau^{-1}$, and the radius and witness carry $\tau$.
No positive lower bound for $\tau$ has been used.
\end{proof}

\begin{corollary}[Uniform logarithmic flow and grid count]
\label{inf:log-flow-count}
For the inputs above, the signed slope is at least $c\tau/p$ in
$|\lambda|<\tau/(16p)$. For any fixed $5/3<q_{\mathrm{gap}}<2$,
all but
\begin{equation}\label{inf:log-bad-splits}
 C\sqrt{1+\tau}\bigl(p^{2/3}+p^{8/3-q_{\mathrm{gap}}}\bigr)
\end{equation}
interior integer splits have actual quartic gap at least
$\kappa_{\mathrm{gap}}\tau p^{-q_{\mathrm{gap}}}$.
\end{corollary}
\begin{proof}
Rows 13--22 of Lemma~\ref{inf:detuning-dependencies} establish the
flow, count, and spacing hypotheses. Apply the clipped-curve
argument of Lemma~\ref{inf:many-splits} with its explicitly stated
number of curves. The expression in \eqref{inf:log-bad-splits}
is $o(p)$ because $q_{\mathrm{gap}}>5/3$.
\end{proof}
